The context effects cannot be learned at this step as learning is individualized—recall that only one of the two contexts is observed for each participant, since $c_i$ is a time-invariant property, such as a demographic variable. 

Therefore, the learning algorithm will assume the same beliefs as it does in the vanilla MAB setting. It will again use Normal-Normal conjugacy to update our beliefs about the distribution of each participant's parameters as we described previously in Section 2.2. 

The learning algorithm will again output an estimate, which we will call $\hat{\theta}_{i,a}^{L}$, and a measure of uncertainty about this estimate, which we will call $\sigma_{i,a,t}^{2,L}$. The action selection strategy therefore still receives two values (estimate and uncertainty) per participant per arm at each time point.

% \begin{figure}[t]
%     \centering
%     \includegraphics[width=0.8\linewidth, trim=0 0 0 0, clip]{final_regret_vs_distance_ratio.png}
%     \caption{Priors were set as $\mu_0 = [0,0], \tau_0 = [1,1]$ with 500 per $\sigma$ value.}
%     \label{fig:ratio}
% \end{figure}

\subsectiontitle{Action Selection Strategy}

The empirical Bayes action selection strategy believes that the estimate output by the learning algorithm, referred to as $\hat{\theta}_{i,a}^{L}$, is an estimate of the mean of $\theta_{i,a}^{L}=\boldsymbol{\theta}_{i,a}^{\top}x_i$. These estimates are distributed around the true mean:
$$\hat{\theta}_{i,a,t}^{L}\mid\theta_{i,a}
\sim
\mathcal{N}(x_i^\top\theta_{i,a},\sigma_{i,a,t}^{2,L})$$

Furthermore, EB believes that $\boldsymbol{\theta}_{i,a} \sim \mathcal{N}(\boldsymbol{\mu}_a, \boldsymbol{\Sigma}_a)$. EB will empirically determine the hyperdistribution through the following likelihood model where $\ell_p(\tau_a^2)$ is the profile log-likelihood:
$$\ell(\boldsymbol{\mu}_a,\tau_a^2)
=
-\frac12
\sum_{i=1}^n
\left[
\log\!\left(2\pi(\tau_a^2+\sigma^2_{i,a,t})\right)
+
\frac{
\left(\hat{\theta}^{L}_{i,a,t}-\mu_{a,0}-\mu_{a,1}c_i\right)^2
}{
\tau_a^2+\sigma^2_{i,a,t}
}
\right].$$
$$\hat{\tau}_a^2
=
\arg\max_{\tau_a^2\ge 0}\;
\ell_p(\tau_a^2), \quad \hat{\boldsymbol{\mu}}_a
=
\left(
\sum_{i=1}^{n}
(\hat{\tau}_a^2+\sigma^2_{i,a,t})^{-1}
\begin{pmatrix}
1 & c_i\\
c_i & c_i^2
\end{pmatrix}
\right)^{-1}
\left(
\sum_{i=1}^{n}
(\hat{\tau}_a^2+\sigma^2_{i,a,t})^{-1}
\begin{pmatrix}
\hat{\theta}^{L}_{i,a,t}\\
c_i\hat{\theta}^{L}_{i,a,t}
\end{pmatrix}
\right)$$

\textit{Note}: Unfortunately, we are unable to get an estimate for the full covariance matrix beyond $\hat{\tau}^2_a$, given that the learning algorithm does not separately estimate the $\boldsymbol{\theta}_{i,a}$. However, for a more precise estimate, one can estimate $\tau^2_a$ separately by $c_i$.

Then, EB will use this hyperdistribution to update each participant's parameters:
$$\boldsymbol{\theta}_{i,a}
\mid
\hat{\theta}^{L}_{i,a,t}
\sim
\mathcal{N}
\left(\tilde{\boldsymbol{\Sigma}}_{i,a,t}
\left(
\boldsymbol{\Sigma}_a^{-1}\boldsymbol{\mu}_a
+
\frac{1}{\sigma^{2,L}_{i,a,t}}
\boldsymbol{x}_i
\hat{\theta}^{L}_{i,a,t}
\right),
\left(
\boldsymbol{\Sigma}_a^{-1}
+
\frac{1}{\sigma^{2,L}_{i,a,t}}
\boldsymbol{x}_i
\boldsymbol{x}_i^{\top}
\right)^{-1}
\right)$$